\documentclass[preprint, 12pt]{elsarticle}
% \documentclass[5p]{elsarticle}
% \usepackage[utf8]{inputenc}
\usepackage[T1]{fontenc}
\usepackage{graphicx}
\usepackage{amsmath}
\usepackage{amsfonts}
\usepackage{amssymb}
\usepackage{booktabs}
\usepackage{array}
\usepackage{multirow}
\usepackage{float}
\usepackage{url}
\usepackage{hyperref}

\usepackage{algorithm}
\usepackage{algpseudocode}

\journal{Nuclear Instruments and Methods in Physics Research}

\begin{document}

\begin{frontmatter}

\title{
    A Combined Spectral Analysis Method for In-Situ Soil Elemental Content Measurement
    }

\author[uta]{Jose A. Cortes\corref{cor1}}
\ead{jose.cortes@uta.edu}
\cortext[cor1]{Corresponding author}

\author[uta]{Andrzej Korzeniowski}
\ead{korzeniowski@uta.edu}

\author[usda]{Galina Yakubova}
\ead{galina.yakubova@usda.gov}

\author[usda]{Aleksandr Kavetskiy}
\ead{aleksandr.kavetskiy@usda.gov}

\author[usda]{H. Allen Torbert}
\ead{allen.torbert@usda.gov}

\affiliation[uta]{organization={University of Texas at Arlington},
                 city={Arlington},
                 postcode={76019},
                 state={TX},
                 country={USA}}

\affiliation[usda]{organization={USDA ARS NSDL},
                  postcode={36832},
                  city={Auburn},
                  state={AL},
                  country={USA}}



\begin{abstract}

Traditional methods for measuring soil carbon are often time-consuming, expensive, and require laboratory analysis. Neutron-gamma analysis (NGA) offers an alternative for rapid (practically in-situ) and non-destructive measurement of soil carbon. This paper describes mathematical methods to measure carbon levels in stimulated spectra of soil samples. Monte-Carlo simulations were carried out to create soil neutron stimulated gamma spectra using MCNP6.2 code. Two categories of analysis are investigated: peak fitting, and edgemember fitting, a third method is introduced that combines the two into a single result. Evaluation of their effectiveness determines that the mixed model performed best overall.

\end{abstract}
\begin{keyword}
soil elemental content \sep gamma-ray spectroscopy 
\sep Pulsed Fast Thermal Neutron Analysis \sep Tagged Neutron Method 
\sep MCNP6.2 simulation \sep peak fitting 
\sep edgemember fitting \sep non-destructive analysis 
\sep environmental monitoring \sep rapid measurement 
\sep agricultural soils
\end{keyword}

\end{frontmatter}

\section{Introduction}

Soil carbon is a key component of soil health and plays a crucial role in the global carbon cycle. Accurate measurement of carbon is essential for measuring soil quality, and its impact on the environment \cite{lal_soil_2018} Traditional methods for measuring carbon are often time-consuming, expensive, and require laboratory analysis \cite{de_ros_advancing_2025}. In situ spectral analysis offers an alternative for rapid and non-destructive measurement of soil carbon \cite{yakubova_tagged_2019}.

The Mobile Inelastic Neutron Scattering system (MINS) developed by the National soil dynamics laboratory USDA Agricultural Research Service (NSDL USDA-ARS) uses pulsed fast thermal neutron analysis (PFTNA) \cite{yakubova_measuring_2025} or the Tagged Neutron Method (TNM) \cite{yakubova_tagged_2019} to measure soil carbon and other elements. Different experimental techniques are implemented in these methods. Regardless, both methods rely on analysis of gamma spectra providing information about the elemental composition of the soil. The analysis methods used to interpret the gamma-ray spectra are Peak Fitting for PFTNA and Endmember Fitting for TNM. This paper presents the combination of these methods in calibration. We simulate the MINS with experimental arrangements of equipment and its spectral acquisition using MCNP6.2 code. As samples, common soil minerals with varied carbon and moisture content. Then we apply different gamma-ray spectroscopy methods, to evaluate and compare their effectiveness in measuring carbon in soil.

\section{Physical Procedure and Simulation}

The MINS system consists of a neutron generator, gamma-ray detectors, and a data acquisition system mounted on a mobile platform. The neutron generator emits fast neutrons that interact with the soil and the environment. In PFTNA, the detector measures during and between pulses of neutrons. This allows for the separate measurement of an inelastic scattering (INS) and thermo nuclear capture (TNC) spectra. In TNM, the neutron generator emits tagged neutrons, and the detectors are synchronized to capture gamma rays emitted from neutron interactions with the soil only.

\subsection{Simulation in MCNP}

\begin{figure}[H]
\centering
\includegraphics[width=90mm,clip,trim=0cm 5.5cm 0cm 4.8cm]{../Figures/Simulation/MCNPGeometry.png}
\caption{Side View Geometry of MINS Architecture}
\label{fig:mcnp_geometry}
\end{figure}

MCNP6.2 \cite{werner_mcnp_2017} was used to simulate gamma-ray spectra resulting from neutron interaction with elements in soil samples. The resulting gamma rays that reach the detector is recorded. Each simulation modeled a soil sample undergoing neutron activation with a base soil mineral and varying concentrations of carbon and moisture content. The geometry was set up to mimic in situ measurement conditions, with a neutron source placed above a soil slab and a detector positioned to capture emitted gamma rays \cite{kavetskiy_neutron-stimulated_2017}. In Figure \ref{fig:mcnp_geometry}, the side view of the MCNP simulation setup is shown, highlighting the arrangement of the neutron source, soil slab, and detector. The neutron generator produces neutrons in solid angle ~0.4 sr directed down to soil (like in TNM modification of MINS). In this case, neutrons  only interact with the soil, causing the emission of gamma rays characteristic of the elements present in the sample. The gamma-ray detectors capture these emissions, and the data acquisition system records the spectra for analysis.

In Figure \ref{fig:mcnp_geometry_top}, the top view of the MCNP simulation setup is presented, providing a clear view of the spatial arrangement of the components, including the shielding around the neutron source and detector, as well as the wheels used for mobility in field applications. 

\begin{figure}[H]
\centering
\includegraphics[width=90mm, clip, trim=0cm 1.2cm 0cm 2cm]{../Figures/Simulation/top.png}
\caption{Cross section of MCNP simulation setup showing wheels and shielding (Top View)}
\label{fig:mcnp_geometry_top}
\end{figure}

Key simulation parameters include:
\begin{itemize}
\item \textbf{Neutron source energy 14.1 MeV:} API120 portable neutron (D-T generator) generator \cite{kavetskiy_energy_2018}
\item \textbf{Soil slab dimensions:} 200 cm × 200 cm × 50 cm \cite{wielopolski_nuclear_2012}
\item \textbf{Detector type:} LaBr detector
\item \textbf{Tally:} F8 (pulse height tally) for gamma spectra
\end{itemize}


\begin{table}[H]
    \centering
    \caption{Architectural Chemical Identifiers}
    \begin{tabular}{lllr}
    \hline
    Material   & Element   & ID         &   Weight \% \\
    \hline
    PE+Pb      & C         & 6000       &     4.2860 \\
                & H         & 1001       &     0.7140 \\
                & Pb        & 82000      &    95.0000 \\
    \hline
    Pb         & Pb        & 82000      &   100.0000 \\
    \hline
    Boric Acid & H         & 1001       &     4.8535 \\
                & B10       & 5010       &     3.4981 \\
                & B11       & 5011       &    13.9923 \\
                & O         & 8016       &    77.6561 \\
    \hline
    Al         & Al        & 13027      &   100.0000 \\
    \hline
    Fe-tube    & C         & 006000.50c &     0.0400 \\
                & Si        & 014000.60c &     0.5000 \\
                & P         & 015031.50c &     0.0230 \\
                & S         & 016000.60c &     0.0150 \\
                & Cr        & 024000.50c &    18.9998 \\
                & Mn        & 025055.60c &     1.0000 \\
                & Fe        & 026000.50c &    70.1723 \\
                & Ni        & 028000.50c &     9.2499 \\
    \hline
    Fe         & Fe        & 26000      &   100.0000 \\
    \hline
    Detector   & Br79      & 35079      &    29.4600 \\
                & Br81      & 35081      &    30.6900 \\
                & La139     & 57139      &    34.8500 \\
                & Ce140     & 58140      &     5.0000 \\
    \hline
    Wheel      & H         & 1001       &    11.8371 \\
                & C         & 6000       &    88.1629 \\
    \hline
    \end{tabular}
    \label{tab:arch_identifiers}
\end{table}

Table \ref{tab:arch_identifiers} lists the architectural chemicals and their identifiers for the soil elements. Table \ref{tab:chemical_identifiers} lists the identifiers for the soil elements.

\begin{table}[H]
    \centering
    \caption{Soil Chemical Identifiers}
    \label{tab:chemical_identifiers}
    \begin{tabular}{lrrrrrrr}
    \hline
    Element         &    C &    H &    O &    Si &    Na &    Al &     K \\
    \hline
    MCNP Identifier & 6000 & 1001 & 8016 & 14000 & 11023 & 13027 & 19000 \\
    \hline
    \end{tabular}
\end{table}

\subsection{Common Soil Minerals}

To investigate the effectiveness of gamma spectra analysis methods for soil carbon measurement, we simulate common soil minerals, including Quartz, Feldspar and Mica \cite{sposito_chemistry_2008}. This data serves as a foundation for applying various analysis techniques.

\begin{table}[H]
\centering
\caption{The elemental weight composition of common soil minerals}
\label{tab:materials}
\begin{tabular}{l @{\extracolsep{\fill}} rrrrrrr}
\hline
 Material   &   C \% &   H \% &   O \% &   Si \% &   Na \% &   Al \% &   K \% \\
\hline
 Carbon     & 100.0 &   0.0 &   0.0 &    0.0 &    0.0 &    0.0 &   0.0 \\
 Water      &   0.0 &  11.2 &  88.8 &    0.0 &    0.0 &    0.0 &   0.0 \\
 Quartz     &   0.0 &   0.0 &  53.3 &   46.7 &    0.0 &    0.0 &   0.0 \\
 Feldspar   &   0.0 &   0.0 &  48.8 &   32.1 &    8.8 &   10.3 &   0.0 \\
 Mica       &   0.0 &   0.5 &  48.2 &   21.2 &    0.0 &   20.3 &   9.8 \\
\hline
\end{tabular}
\end{table}

Table \ref{tab:materials} presents the elemental weight composition of common soil minerals used in the simulations. To measure the effectiveness of analysis methods for carbon measurement, we simulate combinations of soil minerals with varying carbon and water content. The simulations range from 0\% to 15\% carbon (0\%, 1\%, ..., 15\%) and from 0\% to 20\% (0\%, 5\%, ..., 20\%) water content by weight.

\begin{figure}[H]
\centering
\includegraphics[width=90mm]{../Figures/Simulation/soilternary.png}
\caption{Simulated combinations of soil carbon and water levels}
\label{fig:soil_carbon}
\end{figure}

\subsection{Gamma Ray Spectra}

\begin{figure}[H]
\centering
\includegraphics[width=90mm]{../Figures/Simulation/carbonspectrum.png}
\caption{Example Simulated Carbon Spectrum}\label{fig:carbon_spectrum}
\end{figure}

As shown in Figure \ref{fig:carbon_spectrum}, spectra are a representation of the energy distribution of gamma rays captured by a detector. They are typically displayed as a graph with energy bins (or bands) on the x-axis and counts (or intensity) per neutron (in simulation) on the y-axis. The observed spectrum $O(E)$ can be represented by the integral \cite{meng_inter-comparison_2000}:
\begin{equation}
    O(E) = \int_{0}^{\infty} R(E, E_0) \cdot I(E_0) \, dE_0
    \label{eq:integral}
\end{equation}
Here, $I(E_0)$ is the incident spectrum for all possible incident energies, and $R(E, E_0)$ is the detector response function, which accounts for the efficiency of the detector.

This equation can be discretized as:
\begin{equation}
    \begin{bmatrix}
        O_1 \\
        \vdots \\
        O_m
    \end{bmatrix}
    =
    \begin{bmatrix}
        R_{11} & R_{12} & \cdots & R_{1m} \\
        \vdots & \vdots &        & \vdots \\
        R_{m1} & R_{m2} & \cdots & R_{mm}
    \end{bmatrix}
    \times
    \begin{bmatrix}
        I_1 \\
        \vdots \\
        I_m
    \end{bmatrix}
    +
    \begin{bmatrix}
        \epsilon_1 \\
        \vdots \\
        \epsilon_m
    \end{bmatrix}
    \label{eq:discrete}
\end{equation}
where m is the number of energy bins, $\epsilon_i$ ($i=1,\ldots,m$) is noise and $R_{ij}$ is the probability that an incident gamma-ray having an energy falling into bin $i$ will be detected in bin $j$.

\subsection{Gaussian Energy Broadening as Detector Response}

The F8 tally tracks the energy detected by the system, or incident spectra. To model the defficiency in the detectors observed resolution, a response function is applied. The response function used is Gaussian energy broadening (GEB) is a model for detector response. The effects of which can be seen in Figure \ref{fig:GEBvsIncident}. It is applied in simulations by a random sampling of a Gaussian distribution as such:
\begin{equation}
    E' = E + \Delta E
\end{equation}
where $E$ is the true energy of the incoming particle, $\Delta E$ is a random variable sampled from a Gaussian distribution defined by the full width at half maximum (FWHM):

\begin{equation}
    \mathrm{FWHM} = a + b\sqrt{E + cE^2},
\end{equation}
where $E$ is the energy of the particle. The units of $a$, $b$, and $c$ are MeV, MeV$^{1/2}$, and MeV$^{-1}$, respectively. The energy actually scored is sampled from a Gaussian with that FWHM.
The parameters $a$, $b$, and $c$ were determined experimentally as -0.026198 0.059551 -0.037176 respectively \cite{werner_mcnp_2017}.

\begin{figure}[H]
\centering
\includegraphics[width=90mm]{../Figures/Simulation/GEBvsIncident.png}
\caption{Response Function: Gaussian Energy Broadening (GEB)}\label{fig:GEBvsIncident}
\end{figure}

\subsection{Effects on Spectral Readings}

\begin{figure}[H]
\centering
\includegraphics[width=90mm]{../Figures/Simulation/carboneffect.png}
\caption{Carbon effect on Feldspar Spectrum}
\label{fig:soil_carbon_effect}
\end{figure}

The spectral readings obtained from the MCNP simulations provide a detailed distribution of the gamma-ray emissions from the soil samples per neutron, similar to what would be measured in a real-world scenario. Spectrums can be used to identify the presence of specific elements based on their characteristic energy peaks. 

The effect of the carbon on the spectral readings can be seen near Carbons 4.44 MeV gamma peak in Figure \ref{fig:soil_carbon_effect}. As the carbon content increases, the peak becomes more pronounced. Analysis methods can use this information to determine the carbon content of the soil.

As seen in Figure \ref{fig:carbon_gamma}, peaks from Oxygen and Silicon interfere with the carbon peak, caused by gamma emissions from Oxygen and Silicon in the soil. This has an effect on the analysis that needs to be accounted for when determining the carbon content in the soil.

\begin{figure}[H]
\centering
\includegraphics[width=90mm]{../Figures/Simulation/ElemDetectorReadingsCGamma.png}
\caption{Carbon compared to Oxygen and Silicon gamma peaks}
\label{fig:carbon_gamma}
\end{figure}

\section{Data Analysis Methods}

\subsection{Curve Fitting}

\begin{algorithm*}
\caption{Curve Fitting Procedure}
\begin{algorithmic}[1]
    \Require Data points $(x_i, y_i)$, model function $f(x, \text{params})$.
    \State Initialize parameters $\text{params}_0$.
    \State Define loss function: $L(\text{params}) = \sum_i (y_i - f(x_i, \text{params}))^2$.
    \State Use optimization algorithm to find $\text{params}_\text{opt}$ that minimize $L(\text{params})$.
    \State \Return Fitted parameters $\text{params}_\text{opt}$ and fitted curve $f(x, \text{params}_\text{opt})$.
\end{algorithmic}
\end{algorithm*}

The curve fitting procedure begins by collecting data points $(x_i, y_i)$ and selecting a suitable model function\\ 
$f(x, \text{params})$. Initial parameter values $\text{params}_0$ are chosen, and a loss function is defined as the sum of squared residuals: $L(\text{params}) = \sum_i (y_i - f(x_i, \text{params}))^2$. An optimization algorithm is then employed to determine the parameter set $\text{params}_\text{opt}$ that minimizes the loss function. The result is a set of fitted parameters and the corresponding fitted curve $f(x, \text{params}_\text{opt})$. This method is applied with the \texttt{scipy.optimize.curve\_fit} function from the scipy python package \cite{virtanen_scipy_2020} , which utilizes either the Levenberg-Marquardt algorithm \cite{more_levenberg-marquardt_1978} (for unconstrained problems) or the Trust Region Reflective (TRF) algorithm \cite{branch_subspace_1999} (for constrained problems).



\subsection{Peak Baseline Fitting}


\begin{figure}[H]
\centering
% \includegraphics[width=190mm]{../Figures/Analysis/peak_fitting_feldspar_subplots.png}
\includegraphics[width=\linewidth]{../Figures/Analysis/peak_fitting_feldspar_subplots.png}
\caption{Peak Baseline fitting example for feldspar spectrum at 14\% carbon level.}
\label{fig:peak_fitting}
\end{figure}

Peak fitting as shown in Figure \ref{fig:peak_fitting} involves using the least-squares method in identifying and quantifying the baseline and peaks in the spectral data that correspond to specific soil components \cite{gardner_use_2011}. This method is useful for extracting information about the concentration of individual elements in the soil. For effective peak fitting, the data is filtered to focus on the peak area.

The fitting function is defined as:
\begin{equation}
F_f = F_p + F_b
\end{equation}

where $F_p$ is the peak function (e.g., Gaussian) and $F_b$ is the baseline function (e.g., linear or exponential falloff).

\begin{figure}[H]
\centering
\includegraphics[width=90mm]{../Figures/Analysis/carbon_level_vs_predicted_pf.jpg}
\caption{Peak areas are correlated with soil carbon levels over all mineral/carbon/water combinations}
\label{fig:c_to_peak}
\end{figure}

\begin{table}[H]
\centering
\caption{Function parameterizations for peak fitting}
\label{tab:functions}
\begin{tabular}{ll}
\toprule
Function Type & Example Expression \\
\midrule
Peak & \\
% \hline
Gaussian & $a \cdot \exp(-((x - b)^2)/c^2) + d$ \\
\midrule
Baseline & \\
% \hline
Linear & $ax + b$ \\
Exponential Falloff & $a \cdot \exp(-b \cdot x) + c$ \\
\bottomrule
\end{tabular}
\end{table}


After curve fitting, the baseline function is subtracted from the fitted function to isolate the peak, and the area under the peak is calculated using the Simpsons method implemented by \texttt{scipy.integrate.simpson} \cite{cartwright_simpsons_2017} to quantify the concentration of the corresponding element or compound in the soil. As shown in Figure \ref{fig:c_to_peak}, the peak areas are correlated with the carbon levels in the soil samples. The variance in peak area between samples with similar carbon levels may be attributed to other soil components affecting the peak fitting. This includes elements such as oxygen and silicon, the ladder of which, as in Figure \ref{fig:s_to_peak}, is also shown to have a non-negligible correlation with the carbon peak area.

\begin{figure}[H]
\centering
\includegraphics[width=90mm]{../Figures/Analysis/silicon_level_vs_predicted_pf.jpg}
\caption{Silicon levels are correlated with carbon peak areas over all mineral/carbon/water combinations}
\label{fig:s_to_peak}
\end{figure}

\subsection{Spectral Mixture of Endmembers}

Spectral results can be modeled as a mixture of several component spectra, known as endmembers. This approach is useful in scenarios where the observed spectrum is a mixture of distributions. The linear spectral mixture model \cite{adams_classification_1995} can be expressed as:

\begin{equation}
    O(E)_i = \sum_{k=1}^{n} t_{ik} \, O_{ik} + \epsilon_i
    \qquad \text{and} \qquad
    \sum_{k=1}^{n} t_{ik} = 1,
    \label{eq:mixture}
\end{equation}

Where $O(E)_i$ is the intensity in bin $i$ and can be represented as a weighted sum of the intensities from each endmember $O_{ik}$, with coefficients given by the fractions $t_{ik}$. The term $\epsilon_i$ accounts for error.

\subsection{Spectral Mixture Fitting of Endmembers }

The Spectral Mixture Fitting of Endmembers 
leverages the Spectral Mixture of Endmembers model Eq \ref{eq:mixture} to fit known spectral components to the observed spectrum. This was applied by Bishnoi et al. \cite{bishnoi_modeling_2019} using 3 endmembers. The goal is to express the observed spectrum as a linear combination of several known component spectra (endmembers), each representing a specific soil constituent. To do this, we set $t_{ik} = A_k$, where $A_k$ are the coefficients representing the concentration of each component in the mixture.

The combined spectral function is defined as:
\begin{equation}
F_{c_i} = \sum_{k} A_k \cdot O_{k_i} = A_{C} \cdot O_{C_i} + ... + A_{Na} \cdot O_{Na_i}
\label{eq:combined_spectrum}
\end{equation}

where $F_{c_i}$ is the combined spectral function at bin $i$, and $O_{k_i}$ are the observed spectral functions of individual components at bin $i$. In Figure \ref{fig:elemental_component_fitting}, we illustrate the result of the fitting process. Beginning with the spectrum with the highest coefficient, we stack plot the linear combination of the component spectra weighted by their respective coefficients $A_k$. The fitting is performed using the curve fitting procedure outlined earlier, $F_{c}$ does not match its observed spectrum exactly. THis is further discussed in the conclusion.

\begin{figure}[H]
\centering
\includegraphics[width=140mm]{../Figures/Analysis/elemental_linear_combination_feldspar.png}
\caption{Stack plot of edgemember fitting result for feldspar spectrum at 2\% carbon 10\% water level.}
\label{fig:elemental_component_fitting}
\end{figure}

\begin{figure}[H]
\centering
\includegraphics[width=90mm]{../Figures/Analysis/carbon_level_vs_predicted_ca.jpg}
\caption{Carbon Coefficient is correlated with soil carbon levels over all mineral/carbon/water combinations}
\label{fig:component_predictions}
\end{figure}


The carbon coefficient $A_{C}$ is then used for callibration against known carbon concentrations to predict the carbon content in unknown samples. We can see by Figure \ref{fig:component_predictions} that there is a positive correlation between the calculated carbon coefficients from the SMFE method and the actual carbon levels in the soil samples. Unlike in peak fitting, there is a greater spread of predicted carbon coefficients for similar carbon levels. This is because the effect of silicone and oxygen has a greater effect on the methods calculation.

\subsection{Multivariate Calibration Using Singular Value Decomposition}

We employ a multivariate calibration technique. We extract multiple features from the spectral data, such as peak area via Peak Baseline Fitting and the carbon coefficient from edgemember fitting. The data can be represented in matrix form as:

\[
G = 
\begin{bmatrix}
    g_{11} & g_{12} & \cdots & g_{1n} \\
    g_{21} & g_{22} & \cdots & g_{2n} \\
    \vdots & \vdots & \ddots & \vdots \\
    g_{m1} & g_{m2} & \cdots & g_{mn} \\
\end{bmatrix}
\]
for soil sample $m$ and augmented feature $n$ (peak area, edgemember coefficients). A bias term is also included in the model to account for any constant offset in the data.

The relationship between these features and the carbon content can be modeled as:
\[
y = Gx + \epsilon
\]
where $y$ is the vector of observed carbon content, $x$ is the vector of coefficients to be determined, and $\epsilon$ represents the error term.

To calibrate our model, we employ Singular Value Decomposition (SVD) to 
analyze the relationship between the features found in analysis and the carbon content.

A matrix $G$, can be decomposed using SVD as follows:
\begin{equation}
    G = U \Sigma V^T
\end{equation}

Using this decomposition, we can find the SVD solution to least squares, or convex optimization problems.
This allows us to determine the coefficients that best fit our model to the observed data.

To solve for the coefficients $x$ in the equation $Gx = b$, where $b$ is the observed data, we can use the pseudoinverse of $G$:
\begin{equation}
    x = G^+ b
\end{equation}
where $G^+$ is the pseudoinverse of $G$.
The pseudoinverse can be computed from the SVD components as:
\begin{equation}
    G^+ = V \Sigma^+ U^T
\end{equation}
where $\Sigma^+$ is obtained by taking the reciprocal of the non-zero singular values in $\Sigma$ and transposing the resulting matrix.

For example, if $\Sigma$ is a diagonal matrix with singular values $\sigma_1, \sigma_2, \ldots, \sigma_r$ (and possibly zeros for the remaining entries), then $\Sigma^+$ is given by:
\begin{equation}
    \Sigma^+ = 
    \begin{bmatrix}
        \frac{1}{\sigma_1} & 0 & \cdots & 0 \\
        0 & \frac{1}{\sigma_2} & \cdots & 0 \\
        \vdots & \vdots & \ddots & \vdots \\
        0 & 0 & \cdots & \frac{1}{\sigma_r} \\
        0 & 0 & \cdots & 0 \\
        \vdots & \vdots & & \vdots \\
        0 & 0 & \cdots & 0 \\
    \end{bmatrix}
\end{equation}
where the reciprocals are taken only for the nonzero singular values.

In our implementation, we utilize the \texttt{numpy.linalg.svd} function from the \texttt{NumPy} \cite{harris2020array} Python package to perform the singular value decomposition.

The final coefficients obtained from the SVD calibration are then used to predict the carbon content based on the extracted features from introduced spectral data:
\begin{equation}
    \hat{y} = G_{test} x
\end{equation}
where $\hat{y}$ is the predicted carbon content.

\section{Comparison and Results}

\subsection{Metrics for Model Evaluation}

To compare the performance of different models, we split the data randomly between 70\% training and 30\% testing sets. We vary the combination of methods used, then calibrate the results of those methods to the carbon in the testing set using SVD. We use the Mean Squared Error (MSE) between the defined and predicted carbon values.
We define the error for each tested sample $i$ as:
\begin{equation}
    \text{Error}_i = (y_{i} - \hat{y}_{i})^2
\end{equation}
where $y_{i}$ is the simulated carbon content and $\hat{y}_{i}$ is the predicted carbon content for sample $i$ in the testing set.
The Mean Squared Error (MSE) is then calculated as:
\begin{equation}
    \text{MSE} = \frac{1}{N} \sum_{i=1}^{N} \text{Error}_i
\end{equation}
where $N$ is the total number of tested samples.

\subsection{Changing Variables for Multivariate Analysis}

\begin{figure}[H]
    \centering
    \includegraphics[width=90mm]{../Figures/Analysis/edgemembermultivar.jpg}
    \caption{Multivariate calibration improves accuracy}
    \label{fig:edgemember_multivar_results}
\end{figure}

The effects of multivariate analysis can be seen most clearly with the Edgemember Fitting method as in Figure \ref{fig:edgemember_multivar_results}. By including additional variables such as silicon and oxygen coefficients alongside carbon, the model can better account for overlapping spectral features and interferences, and the predictions become more accurate. The most accurate of both methods was the multivariate analysis including carbon, silicon, and oxygen edgemember coefficients and peak areas as variables respectively as seen in Table \ref{tab:multivar_results}. Combining both the peak areas and the component fitting coefficients into a single model improved the accuracy further.

\begin{table}[H]
\centering
\caption{Mean Squared Error (MSE) Comparison of Different Calibration Methods}
\label{tab:multivar_results}
\begin{tabular}{lrrr}
\toprule
Method & MSE \\
\midrule
\textit{Variable Element:} & \textit{C, Si, O} & \textit{C, Si} & \textit{C} \\
\midrule
Combined & 1.98e-06 \\
Edgemember Coeff. & 4.28e-06 & 2.11e-4 & 7.52e-4 \\
Peak Fitting & 1.00e-05 & 1.37e-05 & 1.79e-05\\
\bottomrule
\end{tabular}
\end{table}

\subsection{Conclusion and Discussion}

This study demonstrates the use of multi-variate calibration to enhance soil carbon content predictions by integrating features from various analysis methods. 

\begin{figure}[H]
\centering
\includegraphics[width=90mm]{../Figures/Analysis/elemental_linear_combination_feldspar_RMSE.png}
\caption{Edgemember Fitting (MSRE Error Metric)}
\label{fig:elemental_linear_combination_feldspar_RMSE}
\end{figure}

In Figure \ref{fig:elemental_component_fitting} we observe that the fitted coefficients do not sum exactly to the observed spectrum. Furthermore, the final coefficients are close to their initial values. This is because the curve fitting error is based on MSE minimization, which is more sensitive to larger values found in lower energy regions. As shown in Figure \ref{fig:elemental_linear_combination_feldspar_RMSE} and Table \ref{tab:msre_results} MSRE (Mean Squared Relative Error) improves this. 

\begin{table}[H]
\centering
\caption{Comparison of coefficients using MSRE fitting}
\label{tab:msre_results}
\begin{tabular}{lrr}
\hline
 Element   &   Coefficient &   True \\
\hline
 C         &       0.0156  &       0.0200  \\
 H         &       0.1616  &       0.0156 \\
 O         &       0.4777  &       0.5130  \\
 Si        &       0.1769  &       0.1862  \\
 Na        &       0.0000  &       0.0000  \\
 Al        &       0.1681  &       0.1788  \\
 K         &       0.0000  &       0.0864 \\
\hline
\end{tabular}
\end{table}

As seen in Table \ref{tab:pred_results_vs_RMSE} in the Univariate case, this results in a lower prediction error, but in the Multivariate case, the error of the predictions remains fixed, while the MSE version improves to be better than its counterpart.

\begin{table}[H]
\centering
\caption{Comparison of Different Edgemember Fitting Error Metrics}
\label{tab:pred_results_vs_RMSE}
\begin{tabular}{lrr}
\toprule
Variables & MSE Fitted & MSRE Fitted \\
\midrule
C Si and O & 4.28e-06 & 1.47e-05 \\
C and Si & 2.11e-04 & 2.28e-05 \\
C only & 7.52e-04 & 2.29e-05 \\
\bottomrule
\end{tabular}
\end{table}

Our results demonstrate a significant performance increase, showing that this combined approach yields more accurate predictions of soil carbon content than either method used alone. However, there are many soil types, the soil minerals listed by Sposito \cite{sposito_chemistry_2008} are only one part of the possible makeups that the MINS would encounter in the field. The simulations presented also do not include variance in density. More detailed simulations and spectra to match PFTNA and TNM methods would help to further validate the results presented here.

The calibration step can further include in its calculations the results of various methods. As the dataset grows, more methods of analysis such as convolutional neural networks as presented in Kim et al. \cite{kim_deep_2025}. Other models, such as deep learning models, or other deconvolution methods can be ensembled together to improve accuracy further.

\bibliographystyle{elsarticle-num}
\bibliography{references}

\end{document}